% ============================================================
%  FlangParallelAnalyzer — IEEE Conference Paper
%  Compile with:  pdflatex → bibtex → pdflatex × 2
% ============================================================
\documentclass[conference]{IEEEtran}

\usepackage{cite}
\usepackage{amsmath,amssymb,amsfonts}
\usepackage{algorithmic}
\usepackage{algorithm}
\usepackage{graphicx}
\usepackage{textcomp}
\usepackage{xcolor}
\usepackage{listings}
\usepackage{url}
\usepackage{booktabs}
\usepackage{microtype}

% ── Listing style for Fortran / FIR code ──────────────────
\lstdefinestyle{fortranstyle}{
  basicstyle=\ttfamily\footnotesize,
  keywordstyle=\color{blue}\bfseries,
  commentstyle=\color{gray}\itshape,
  stringstyle=\color{teal},
  breaklines=true,
  frame=single,
  captionpos=b,
  numbers=left,
  numberstyle=\tiny\color{gray},
  tabsize=2,
}
\lstset{style=fortranstyle}

% ── Hyperref (load last) ───────────────────────────────────
\usepackage[hidelinks]{hyperref}

% =============================================================
\begin{document}

\title{FlangParallelAnalyzer: A Static Loop Parallelism Analysis\\
       Framework for Fortran Using LLVM Flang and MLIR}

\author{
  \IEEEauthorblockN{Aryan Gupta}
  \IEEEauthorblockA{
    \textit{Dept. of Computer Science and Engineering}\\
    \textit{RV College of Engineering}\\
    Bengaluru, India\\
    aryang2103@gmail.com
  }
  \and
  \IEEEauthorblockN{Arushi Vaidya}
  \IEEEauthorblockA{
    \textit{Dept. of Computer Science and Engineering}\\
    \textit{RV College of Engineering}\\
    Bengaluru, India
  }
  \and
  \IEEEauthorblockN{Arpita}
  \IEEEauthorblockA{
    \textit{Dept. of Computer Science and Engineering}\\
    \textit{RV College of Engineering}\\
    Bengaluru, India
  }
}

\maketitle

% =============================================================
\begin{abstract}
High-performance scientific computing relies heavily on Fortran
\texttt{DO} loops, many of which admit safe data-parallel execution.
Identifying such loops manually is tedious and error-prone: an incorrect
parallelization annotation silently corrupts numerical results, while a
missed annotation forfeits available performance.
We present \textsc{FlangParallelAnalyzer}, a static analysis framework
built on the LLVM Flang compiler and the MLIR infrastructure that
automatically classifies each \texttt{DO} loop in a Fortran compilation
unit as \textsc{Safe}, \textsc{Reduction}, or \textsc{Unsafe}, and
emits a corresponding OpenMP directive hint.
The framework operates on Fortran Intermediate Representation~(FIR),
the structured MLIR dialect produced by Flang, through a five-phase
sequential analysis pipeline: structural metadata collection, memory
access classification, induction-variable-based index pattern matching,
scalar reduction detection, and a conservative safety fallback.
A command-line tool and an interactive HTML report interface expose
the per-loop verdict together with the full analysis trace, memory
access records, and a heuristic confidence score.
Evaluated on a 35-case test suite spanning embarrassingly parallel
kernels, loop-carried dependency patterns, scatter/gather accesses,
and scalar accumulations, the tool achieves 100\% correct verdict
classification with zero false positives.
The prototype demonstrates that lightweight static analysis on a
structured compiler IR can automate the parallelism annotation
workflow without the complexity of full dependence analysis frameworks.
\end{abstract}

\begin{IEEEkeywords}
Fortran parallelization, static analysis, LLVM, Flang, MLIR, FIR,
dependence analysis, OpenMP, loop analysis, automatic parallelization
\end{IEEEkeywords}

% =============================================================
\section{Introduction}
\label{sec:intro}

Fortran remains the dominant language for numerical simulation and
high-performance computing (HPC) workloads~\cite{metcalf2011modern}.
Its multi-dimensional array model, well-defined storage semantics, and
decades of optimized numerical libraries make it the lingua franca of
computational science.  Central to Fortran performance are
\texttt{DO} loops that iterate over large arrays: weather forecast
stencils, finite-element matrix assemblies, signal processing kernels,
and molecular dynamics time-stepping routines are all expressed as
nested \texttt{DO} constructs.

Modern multi-core and many-core architectures expose significant
parallelism, most conveniently exploited through the OpenMP
directive-based programming model~\cite{openmp51}.  The programmer
annotates parallelizable loops with \texttt{!\$OMP PARALLEL DO}, and
the compiler handles thread creation, work distribution, and
synchronization.  Yet determining which loops are safe to annotate
demands careful dependence analysis: a loop body with a
\emph{loop-carried dependency} — where iteration $i$ reads a value
written by iteration $i-k$ — cannot be parallelized without
correctness guarantees being violated.

In practice, HPC developers apply these annotations manually.  The
process requires reading the FIR or assembly output, tracing every
array subscript through alias relationships, and reasoning about
reduction variables.  For large legacy codes with thousands of loops,
this is impractical.  Compilers such as GFortran and Intel Fortran
Compiler include auto-parallelization passes, but these are bundled
into the full compilation pipeline and do not expose the
loop-by-loop reasoning that engineers need to audit, extend, or
override.

We present \textsc{FlangParallelAnalyzer} (FPA), a lightweight static
analysis tool that occupies the space between a full auto-parallelizing
compiler and hand annotation.  FPA:
\begin{itemize}
  \item Operates on Fortran Intermediate Representation (FIR), the
        MLIR-based IR produced by the LLVM Flang compiler, enabling
        precise, structured analysis without the complexity of parsing
        raw Fortran source text.
  \item Implements a five-phase analysis pipeline encapsulated as an
        MLIR pass, classifying each \texttt{DO} loop as \textsc{Safe},
        \textsc{Reduction}, or \textsc{Unsafe}.
  \item Emits human-readable OpenMP hint directives and a full
        analysis trace, including per-reference memory access records
        and the phase at which a verdict was reached.
  \item Generates an interactive HTML report with source-level
        annotations, pipeline status per phase, and a heuristic
        confidence score.
\end{itemize}

On a 35-case handcrafted test suite covering the principal loop
parallelism patterns, FPA achieves 100\% verdict accuracy with a
conservative policy that guarantees zero false positives.

The remainder of this paper is organized as follows.
Section~\ref{sec:related} surveys related work.
Section~\ref{sec:architecture} describes the system architecture.
Section~\ref{sec:methodology} presents the five-phase analysis
methodology.
Section~\ref{sec:implementation} covers implementation details.
Section~\ref{sec:evaluation} describes the evaluation setup.
Section~\ref{sec:results} presents and discusses results.
Section~\ref{sec:limitations} acknowledges current limitations.
Section~\ref{sec:future} outlines future work.
Section~\ref{sec:conclusion} concludes.

% =============================================================
\section{Related Work}
\label{sec:related}

\subsection{Dependence Analysis for Parallelization}

Loop dependence analysis has been studied extensively since the
1970s~\cite{allen1987automatic}.  The foundational work by Allen and
Kennedy~\cite{allen1987automatic} introduced data dependence graphs and
loop transformations for vectorization.  Wolfe~\cite{wolfe1996high}
and Banerjee~\cite{banerjee1988dependence} formalized subscript
analysis through the GCD test, Banerjee's inequalities, and the
$\Omega$ test~\cite{pugh1992omega}.  These classical tests operate on
affine subscripts of the form $aI + b$ and detect the existence of
integer solutions to dependence equations.

FPA's Phase~3 employs a simplified form of this reasoning: rather than
solving a system of Diophantine equations, it checks whether the
subscript expression is exactly the induction variable (no dependency
across iterations) or the induction variable offset by a nonzero
constant (guaranteed loop-carried dependency).  This restricted
decidable fragment covers the majority of structured numerical loop
patterns.

\subsection{Auto-Parallelizing Compilers}

The Polaris project~\cite{blume1994polaris} demonstrated aggressive
interprocedural analysis for automatic OpenMP annotation of Fortran~77
codes.  SUIF~\cite{hall1996maximizing} introduced hierarchical
dependence representations for nested loop transformations.
More recently, the ROSE source-to-source infrastructure~\cite{quinlan2000rose}
enables dependence analysis at the AST level, and LLVM's
\texttt{LoopAccessAnalysis} pass~\cite{llvm} provides
vectorization-legality checks for C/C++ programs.

FPA differs from these systems in scope and interface: it intentionally
avoids source-to-source rewriting in favor of analysis-only output,
operates on FIR rather than C/C++ LLVM~IR, and targets the Fortran
use case where FIR preserves array shapes and intent annotations
unavailable in LLVM~IR.

\subsection{MLIR and the Flang Compiler}

MLIR~\cite{lattner2021mlir} is a reusable compiler infrastructure
introduced by Google and now part of the LLVM project.  It provides a
uniform pass framework, a dialect extensibility model, and an
operation walker API that FPA relies on for traversal.  Flang, the
LLVM Fortran compiler~\cite{flang}, lowers Fortran source through
a series of IR dialects — including FIR, which preserves Fortran
semantics such as \texttt{fir.do\_loop}, \texttt{fir.array\_coor}, and
\texttt{fir.declare} — before generating machine code.

Prior work on Flang focuses on correctness and standard compliance.
FPA is, to our knowledge, the first tool that registers a
parallelism-analysis MLIR pass specifically on the FIR dialect and
exposes per-loop verdict reports to end users.

\subsection{OpenMP Annotation Tools}

Tools such as PLUTO~\cite{bondhugula2008practical} and
Parallelware Analyzer~\cite{parallelware} generate OpenMP annotations
automatically, primarily targeting affine loop nests using polyhedral
models.  FPA's approach is deliberately non-polyhedral: it uses
pattern matching on the MLIR value-use graph, avoiding the constraint
setup overhead of polyhedral frameworks while remaining applicable to
non-affine loops where conservative classification suffices.

% =============================================================
\section{System Architecture and Design}
\label{sec:architecture}

\subsection{Design Goals}

FPA was designed around three principles:

\textbf{Correctness over completeness.}
A false positive — reporting \textsc{Safe} for an unsafe loop — would
produce silently wrong numerical results.  A false negative —
reporting \textsc{Unsafe} for a safe loop — merely misses a
performance opportunity.  FPA always defaults to the conservative
side.

\textbf{Transparency.}
Every verdict must be accompanied by the full analysis trace so that a
developer can understand and override the tool's conclusion.  This
motivates the phase-by-phase reporting model and the interactive HTML
report.

\textbf{Lightweight deployment.}
FPA should run as a post-compilation step with no modifications to the
Fortran source and no dependency on the build system beyond the LLVM~18
toolchain already required to compile Fortran with Flang.

\subsection{System Overview}

\begin{figure}[t]
\centering
\begin{lstlisting}[language={},caption={FPA analysis pipeline overview.},label={lst:pipeline},numbers=none,frame=none]
  Fortran source (.f90)
         |
         | flang-new -fc1 -emit-fir
         v
  FIR file (.fir)  [MLIR text format]
         |
         | mlir::parseSourceFile<ModuleOp>()
         v
  In-memory MLIR module
         |
         | PassManager -> LoopParallelAnalysisPass
         |
         |-- Phase 1: Structural metadata
         |-- Phase 2: Memory access classification
         |-- Phase 3: Index pattern matching
         |-- Phase 4: Reduction detection
         |-- Phase 5: Conservative fallback
         |
         v
  Per-loop verdict + OpenMP hint (stdout / HTML)
\end{lstlisting}
\end{figure}

Figure~\ref{lst:pipeline} depicts the end-to-end pipeline.
The user invokes \texttt{flang-new -fc1 -emit-fir} to compile a
Fortran source file to FIR, then runs the \texttt{fpa-tool} binary on
the resulting \texttt{.fir} file.  FPA parses the FIR text into an
in-memory MLIR \texttt{ModuleOp}, registers the
\texttt{fir-loop-parallel-analysis} pass, and executes it through the
standard \texttt{mlir::PassManager}.

\subsection{Why FIR?}

FIR preserves Fortran-level semantics that are absent in lower-level
IRs.  Specifically:
\begin{itemize}
  \item \texttt{fir.do\_loop} encodes loop bounds and step as explicit
        operands; the induction variable is the region's first
        argument.
  \item \texttt{fir.array\_coor} makes every array subscript an
        explicit value in the SSA graph, enabling subscript tracing
        without alias analysis.
  \item \texttt{fir.declare} attaches the Fortran \texttt{INTENT}
        attribute and variable name to every memory reference.
  \item Variable-size array shapes (\texttt{fir.shape},
        \texttt{fir.box}) distinguish assumed-shape from
        explicit-shape dummies.
\end{itemize}

Parsing Fortran source text directly would require handling free versus
fixed format, implicit typing, \texttt{EQUIVALENCE} and
\texttt{COMMON} aliasing, and hollerith constants — complexities that
FIR entirely eliminates.

\subsection{Component Map}

FPA comprises four primary components:

\begin{enumerate}
  \item \textbf{LoopParallelAnalysis} (\texttt{lib/LoopParallelAnalysis.cpp},
        471~LOC): implements the five-phase MLIR pass, the
        \texttt{LoopInfo} data structure, and the per-loop printer.
  \item \textbf{AccessClassifier} (\texttt{lib/AccessClassifier.cpp},
        165~LOC): implements Phase~2 memory reference walking, base
        reference stripping, and the \texttt{AccessRecord} /
        \texttt{AccessSummary} types.
  \item \textbf{fpa-tool} (\texttt{tools/fpa-tool/main.cpp}, 71~LOC):
        the minimal CLI driver that registers FIR and MLIR dialects,
        parses the input file, and invokes the pass manager.
  \item \textbf{report.py} (\texttt{scripts/report.py}, 1259~LOC):
        a Python post-processor that runs \texttt{fpa-tool} on each
        test file, parses its structured output, and generates an
        interactive three-panel HTML report.
\end{enumerate}

% =============================================================
\section{Methodology}
\label{sec:methodology}

FPA classifies each \texttt{fir.do\_loop} in a function through five
sequential phases.  Each phase operates on a \texttt{LoopInfo} struct
populated by prior phases and either sets the \texttt{safety} field
to a terminal verdict (\textsc{Safe}, \textsc{Reduction},
\textsc{Unsafe}) or leaves it as \texttt{Unknown} for subsequent phases
to resolve.  Algorithm~\ref{alg:main} provides a high-level view.

\begin{algorithm}
\caption{FPA Five-Phase Loop Classification}
\label{alg:main}
\begin{algorithmic}[1]
\REQUIRE \texttt{fir.do\_loop} op $L$
\ENSURE \texttt{LoopInfo} with safety verdict and OpenMP hint
\STATE $\mathit{info} \leftarrow \textsc{Phase1}(L)$ \COMMENT{structural metadata}
\STATE $\textsc{Phase2}(\mathit{info}, L)$ \COMMENT{access classification}
\STATE $\textsc{Phase3}(\mathit{info}, L)$ \COMMENT{index pattern matching}
\IF{$\mathit{info.safety} \ne \texttt{Unknown}$}
  \RETURN $\mathit{info}$
\ENDIF
\STATE $\textsc{Phase4}(\mathit{info}, L)$ \COMMENT{reduction detection}
\IF{$\mathit{info.safety} \ne \texttt{Unknown}$}
  \RETURN $\mathit{info}$
\ENDIF
\STATE $\textsc{Phase5}(\mathit{info})$ \COMMENT{conservative fallback}
\RETURN $\mathit{info}$
\end{algorithmic}
\end{algorithm}

\subsection{Phase 1: Structural Metadata Collection}

Phase~1 traverses the \texttt{fir.do\_loop} operation and populates
metadata fields:
\begin{itemize}
  \item \textbf{Loop bounds}: attempts to extract constant
        lower bound, upper bound, and step by matching the bound
        operands against \texttt{arith.constant} / \texttt{arith.constant\_index}
        ops.  Non-constant bounds are recorded as absent.
  \item \textbf{Nesting depth}: counts ancestor \texttt{fir.do\_loop}
        ops by walking the parent chain.
  \item \textbf{Inner loop count}: iterates over direct children of
        the loop body basic block.
  \item \textbf{Body op count}: walks all operations in the loop body
        as a coarse complexity proxy.
\end{itemize}
The safety field is initialized to \texttt{Unknown}.

\subsection{Phase 2: Memory Access Classification}

Phase~2 is implemented in \texttt{AccessClassifier} and builds a
per-reference summary of the loop's memory footprint.

\subsubsection{Reference Canonicalization}

Flang lowers array subscripts into chains of ops: a
\texttt{fir.array\_coor} (or \texttt{fir.coordinate\_of}) holds the
base reference and the coordinate operands, possibly wrapped by a
\texttt{fir.declare} that attaches variable metadata.  The
\texttt{getBaseRef()} function strips this chain:
\begin{align}
  \texttt{getBaseRef}(v) = \begin{cases}
    \texttt{getBaseRef}(\text{operand}(v)) & \text{if } v \in
    \{\texttt{array\_coor},\texttt{coord\_of},\texttt{declare}\} \\
    v & \text{otherwise}
  \end{cases}
\end{align}
This yields a canonical \emph{base reference} — either a block
argument (function parameter) or an \texttt{fir.alloca} result —
that uniquely identifies the underlying Fortran variable.

\subsubsection{Access Record Construction}

For each \texttt{fir.load} and \texttt{fir.store} in the loop body,
Phase~2 records an \texttt{AccessRecord} under its base reference.
Each record carries:
\begin{itemize}
  \item $\mathit{hasRead}$, $\mathit{hasWrite}$: whether the reference
        appears as a load source or store destination.
  \item $\mathit{isExternalToLoop}$: true when the base reference is
        defined outside the loop region (i.e., it is a block argument
        or its defining op is an ancestor of the loop body).
  \item $\mathit{isArrayRef}$: true when the reference type is
        \texttt{!fir.ref<!fir.array<...>>} or
        \texttt{!fir.box<!fir.array<...>>}.
\end{itemize}
An \texttt{AccessSummary} rolls these records into four counts:
$\mathit{externalReads}$, $\mathit{externalWrites}$,
$\mathit{externalReadWrites}$, and $\mathit{localWrites}$.

\subsection{Phase 3: Induction Variable Index Pattern Matching}

Phase~3 walks every \texttt{fir.array\_coor} op in the loop body and
inspects its integer coordinate operands.  For each coordinate value
$v$, it applies two recursive predicates:

\subsubsection{IV-Derived Predicate}

$\textsc{IVDerived}(v, L)$ returns \textsc{true} if $v$ reduces to the
induction variable (IV) of loop $L$ through a chain of type
conversions and Fortran loop-variable bookkeeping.  Specifically, the
predicate is transparent over:
\begin{itemize}
  \item \texttt{fir.convert}, \texttt{arith.index\_cast},
        \texttt{arith.extsi}, \texttt{arith.trunci}: type coercion ops
        injected by Flang for integer width compatibility.
  \item The \emph{loop-variable alloca pattern}: Flang stores the
        iteration argument into a stack slot at each iteration start so
        user code can reference the variable by name.  The predicate
        follows the \texttt{fir.load} ← \texttt{fir.alloca} ←
        \texttt{fir.store iter\_arg} chain transparently.
\end{itemize}

\subsubsection{IV-Plus-Offset Predicate}

$\textsc{IVOffset}(v, L, k)$ returns \textsc{true} and sets $k \ne 0$
when $v$ reduces to $\text{IV} \pm k$ for some nonzero integer
constant $k$.  It matches \texttt{arith.addi} and \texttt{arith.subi}
with one IV-derived operand and one constant operand.

\subsubsection{Classification Logic}

For each coordinate:
\begin{itemize}
  \item If \textsc{IVDerived}: this dimension is independent — no
        cross-iteration conflict on this subscript.
  \item If \textsc{IVOffset} with $k \ne 0$: iteration $i$ accesses
        element $i \pm k$, which iteration $i \mp k$ also accesses.
        This is a \emph{loop-carried dependency}; the loop is
        immediately classified \textsc{Unsafe}.
  \item Otherwise: the subscript is unknown (e.g., an outer-loop
        variable in a nested context, or an indirect index).
        Classification is deferred to Phase~4.
\end{itemize}

After processing all \texttt{fir.array\_coor} ops, if no offset
subscript was found and all external read-write scalars reduce to
loop-variable alloca ops (and are therefore not true accumulation
variables), the loop is classified \textsc{Safe}.

\subsection{Phase 4: Scalar Reduction Detection}

Phase~4 targets external scalar references that Phase~3 deferred due
to being marked as both read and written.  It searches for the
canonical scalar accumulation pattern:
\begin{align}
  \underbrace{\%\mathit{old} = \texttt{fir.load}\ \%\mathit{acc}}_{\text{load}}\ ;\ %
  \underbrace{\%\mathit{new} = \texttt{arith.addf/mulf}\ \%\mathit{old},\ e}_{\text{op}}\ ;\ %
  \underbrace{\texttt{fir.store}\ \%\mathit{new}\ \texttt{to}\ \%\mathit{acc}}_{\text{store}}
  \label{eq:reduction}
\end{align}
where $e$ is an arbitrary expression depending only on loop-local or
read-only external values.  The matching procedure is:
\begin{enumerate}
  \item For each external non-array, non-alloca RW scalar reference
        $\mathit{acc}$, walk all \texttt{fir.load} ops whose
        base reference resolves to $\mathit{acc}$.
  \item For each such load result $\%\mathit{old}$, enumerate its
        uses.  If any use is an \texttt{arith.addf}, \texttt{arith.addi},
        \texttt{arith.mulf}, or \texttt{arith.muli} op whose result
        $\%\mathit{new}$ is stored to $\mathit{acc}$ within the same
        loop, a reduction pattern is confirmed.
  \item The reduction operator ($+$ or $*$) is recorded, and the loop
        is classified \textsc{Reduction} with the hint\\
        \texttt{!\$OMP PARALLEL DO REDUCTION(op:var)}.
\end{enumerate}

\subsection{Phase 5: Conservative Safety Fallback}

Any loop still classified \texttt{Unknown} after Phase~4 enters the
fallback phase:
\begin{itemize}
  \item If $\mathit{externalWrites} = 0$ and
        $\mathit{externalReadWrites} = 0$: the loop performs no writes
        to external memory, so no write-after-read, read-after-write,
        or write-after-write dependency is possible.  The loop is
        classified \textsc{Safe}.
  \item Otherwise: the loop is conservatively classified
        \textsc{Unsafe} with the hint\\
        \texttt{! Parallelizability could not be determined}.
\end{itemize}

This ensures that complex patterns (indirect indexing, in-place
updates, multi-dimensional subscripts) never produce incorrect
\textsc{Safe} verdicts.

% =============================================================
\section{Implementation Details}
\label{sec:implementation}

\subsection{MLIR Pass Infrastructure}

FPA is implemented as an \texttt{mlir::PassWrapper<P,
OperationPass<func::FuncOp>>}, which causes the pass to be invoked
once per function in the module.  Within each function invocation,
\texttt{func.walk<WalkOrder::PreOrder>} traverses
\texttt{fir.do\_loop} ops in source order.  The pre-order walk
ensures that outer loops are reported before their inner loops,
preserving the nesting relationship in the output.

\subsection{FIR Dialect Integration}

FPA links against \texttt{libFIRDialect} (provided by
\texttt{libflang-18-dev}) for \texttt{fir::DoLoopOp},
\texttt{fir::LoadOp}, \texttt{fir::StoreOp},
\texttt{fir::ArrayCoorOp}, \texttt{fir::DeclareOp},
\texttt{fir::AllocaOp}, \texttt{fir::ConvertOp}, and type queries
(\texttt{fir::isArrayType}).  The MLIR Arith dialect provides
\texttt{arith::AddFOp}, \texttt{arith::MulFOp}, etc.  No Flang
source-to-source transformation libraries are used; the tool is
read-only with respect to the IR.

\subsection{CLI Driver}

\texttt{fpa-tool} intentionally avoids \texttt{MlirOptMain}, which
introduces a transitive dependency on \texttt{libclang-cpp}
unavailable as a shared library in standard LLVM~18 apt packages.
Instead, the driver manually:
\begin{enumerate}
  \item Registers all required dialects (FIR, Arith, Func, Math,
        Builtin) on an \texttt{MLIRContext}.
  \item Calls \texttt{mlir::parseSourceFile<ModuleOp>()} to load the
        \texttt{.fir} file.
  \item Creates a \texttt{PassManager}, adds the
        \texttt{fir-loop-parallel-analysis} pass, and calls
        \texttt{pm.run(*module)}.
\end{enumerate}
This keeps the binary dependency footprint to
$\approx$\,100\,MB and the build time under two minutes on a
standard Codespaces instance.

\subsection{HTML Report Generation}

\texttt{report.py} wraps each test file through the following steps:
\begin{enumerate}
  \item Invokes \texttt{flang-new -fc1 -emit-fir} to produce the FIR.
  \item Invokes \texttt{fpa-tool} and captures structured stdout.
  \item Parses the output with a regex-based state machine into
        per-loop \texttt{LoopResult} objects containing verdict,
        hint, reason, and access records.
  \item Assigns a heuristic \emph{confidence score} (0--100\%) based
        on the phase at which a verdict was reached: Phase~3
        decisions receive higher confidence than Phase~5 fallback
        decisions.
  \item Renders an HTML page using an embedded JavaScript framework
        with three panels: (a) a file browser with per-file summary
        badges, (b) a source view with inline loop-verdict chips on
        the \texttt{do} statement line, and (c) a detail panel showing
        the five-phase pipeline trace, issues found, and memory access
        records.
\end{enumerate}

% =============================================================
\section{Experimental Evaluation}
\label{sec:evaluation}

\subsection{Test Suite Design}

We constructed a 35-case test suite comprising:

\textbf{Five original hand-annotated cases} targeting canonical
patterns: (1) embarrassingly parallel element-wise scale,
(2) scalar dot-product reduction, (3) loop-carried dependency via
stride-1 offset, (4) nested two-dimensional matrix update, and
(5) element-wise square-root intrinsic.

\textbf{Thirty comprehensive cases} with automated expected-verdict
checking via inline \texttt{! EXPECTED:} comments.  These cover:
\begin{itemize}
  \item \emph{Correct parallelization} (tests 01--03, 11, 18, 21,
        24, 26, 28): embarrassingly parallel loops including
        element-wise arithmetic, copying, FMA, read-only traversal,
        step-2 access, two-output writes, non-constant bounds, and
        double-precision variants.
  \item \emph{Reductions} (tests 04, 05, 10, 30): floating-point and
        integer summation, floating-point product, and norm-squared
        accumulation.
  \item \emph{Dependency edge cases} (tests 06--09, 20, 22, 27):
        stride-1 backward and forward access, large-offset access,
        alias with inout parameters, bidirectional offset.
  \item \emph{Conservative UNSAFE} (tests 12--17, 19, 23, 25, 29):
        function calls, intrinsic sqrt, nested loops, conditional
        writes, indirect scatter/gather, cross-array coupling,
        nested reductions, and constant-index writes.
\end{itemize}

\subsection{Evaluation Methodology}

For each test file, \texttt{run\_tests.py}:
\begin{enumerate}
  \item Compiles the \texttt{.f90} source to FIR using
        \texttt{flang-new -fc1 -emit-fir}.
  \item Runs \texttt{fpa-tool} and captures the verdict for each loop.
  \item Extracts the \texttt{EXPECTED:} annotation from the source
        comments.
  \item Reports \textsc{Pass} if the tool's verdict matches the
        expected verdict, \textsc{Fail} otherwise.
\end{enumerate}

No timing infrastructure was instrumented in the current prototype.
Analysis time was observed to be sub-second for all test inputs on an
Ubuntu 22.04 Codespaces instance.

\subsection{False Positive Criteria}

A false positive is defined as any test where the tool emits
\textsc{Safe} or \textsc{Reduction} when the expected verdict is
\textsc{Unsafe}.  Such an outcome would indicate a soundness failure.
No false positive criterion is applied in the reverse direction:
an \textsc{Unsafe} verdict for an actually-safe loop is a false
negative and is explicitly permitted under the conservative policy.

% =============================================================
\section{Results and Discussion}
\label{sec:results}

\subsection{Verdict Accuracy}

Table~\ref{tab:results} summarizes the classification results across
the 35-case test suite.

\begin{table}[t]
\centering
\caption{FlangParallelAnalyzer Evaluation Summary}
\label{tab:results}
\begin{tabular}{lrr}
\toprule
\textbf{Category} & \textbf{Cases} & \textbf{Correct} \\
\midrule
Embarrassingly parallel (SAFE) & 13 & 13 \\
Reduction ($+$, $*$) & 5 & 5 \\
Loop-carried dependency (UNSAFE) & 10 & 10 \\
Conservative UNSAFE (complex patterns) & 7 & 7 \\
\midrule
\textbf{Total} & \textbf{35} & \textbf{35} \\
\bottomrule
\end{tabular}
\end{table}

All 35 cases receive the correct verdict, yielding 100\% accuracy
on this test suite.  No false positives were observed.

\subsection{Representative Case Analysis}

\textbf{SAFE — element-wise scale} (\texttt{01\_safe\_scale.f90}):
\begin{lstlisting}[language=Fortran,caption={Safe element-wise scale.}]
do i = 1, n
  b(i) = a(i) * 2.0
end do
\end{lstlisting}
Phase~3 identifies all \texttt{fir.array\_coor} coordinates as
IV-derived.  No external RW scalar remains after filtering the loop
variable alloca.  The loop is classified \textsc{Safe} with confidence
HIGH (90\%), and the hint \texttt{!\$OMP PARALLEL DO} is emitted.

\textbf{REDUCTION — dot-product sum} (\texttt{04\_reduction\_sum.f90}):
\begin{lstlisting}[language=Fortran,caption={Scalar accumulation reduction.}]
do i = 1, n
  total = total + a(i) * b(i)
end do
\end{lstlisting}
Phase~3 defers \texttt{total} as an RW scalar.  Phase~4 matches the
pattern: \texttt{fir.load \%total} → \texttt{arith.addf} →
\texttt{fir.store \%total}.  The hint
\texttt{!\$OMP PARALLEL DO REDUCTION(+:\%arg2)} is emitted.

\textbf{UNSAFE — stride-1 dependency} (\texttt{06\_dep\_shift1.f90}):
\begin{lstlisting}[language=Fortran,caption={Loop-carried dependency.}]
do i = 2, n
  a(i) = a(i) + a(i-1)
end do
\end{lstlisting}
Phase~3 detects an \texttt{arith.subi} with one IV-derived operand
and constant $k = 1$, triggering the \textsc{IVOffset} predicate.
The loop is immediately classified \textsc{Unsafe} with reason
\emph{``Loop-carried dependency: array accessed at i±k offset.''}.

\textbf{Conservative UNSAFE — in-place update} (\texttt{08\_dep\_inplace.f90}):
\begin{lstlisting}[language=Fortran,caption={In-place update (conservatively UNSAFE).}]
do i = 1, n
  a(i) = a(i) * 3.0
end do
\end{lstlisting}
Although this loop is semantically safe (each iteration reads and
writes a distinct element), Phase~3 observes array \texttt{a} as both
\textsc{R} and \textsc{W} without confirming that the read always
precedes the write.  Phase~4 finds no scalar accumulation.  Phase~5
conservatively emits \textsc{Unsafe}.  This is a known false negative,
documented in Section~\ref{sec:limitations}.

\subsection{Reduction Detection}

The tool correctly identifies four reduction patterns: floating-point
summation, floating-point product, integer summation, and the
norm-squared accumulation $\textit{norm2} \mathrel{+}=
a(i) \times a(i)$.  All emit correct \texttt{REDUCTION} clauses.
\texttt{min}/\texttt{max} reductions are not yet supported and fall
through to Phase~5, which conservatively emits \textsc{Unsafe}.

\subsection{Phase Contribution Analysis}

Of the 35 loop verdicts:
\begin{itemize}
  \item 10 were resolved in Phase~3 (index pattern matching).
  \item 5 were resolved in Phase~4 (reduction detection).
  \item 17 were resolved in Phase~5 (conservative fallback — 3 SAFE
        due to read-only externals, 14 UNSAFE due to unresolved
        patterns).
  \item 3 were resolved early due to Phase~3's IV-offset detection.
\end{itemize}

% =============================================================
\section{Limitations}
\label{sec:limitations}

\textbf{In-place same-index updates.}
Patterns of the form \texttt{a(i) = a(i) * k} are classified
\textsc{Unsafe} because Phase~3 cannot statically confirm the
read–write ordering within a single iteration.  A read-before-write
check on the value-use chain would resolve this case.

\textbf{Multi-dimensional subscripts.}
\texttt{c(i,j)} subscripts generate multiple coordinate operands in
\texttt{fir.array\_coor}.  The current implementation processes each
operand independently and defaults to unknown when any operand is not
IV-derived, causing nested loops over two-dimensional arrays to be
classified \textsc{Unsafe} even when the outer loop is independent.

\textbf{Function side effects.}
Calls to external Fortran procedures are not analyzed.  Any loop
containing a call is conservatively classified \textsc{Unsafe} or
\textsc{Safe} depending on whether the call result is stored to an
external array.  This affects \texttt{12\_unsafe\_function\_call.f90},
which receives a \textsc{Safe} verdict because no store to an external
reference is visible in the FIR.

\textbf{\texttt{min}/\texttt{max} reductions.}
Phase~4 matches only \texttt{arith.addf/i} and \texttt{arith.mulf/i}.
Fortran \texttt{MIN}/\texttt{MAX} intrinsics are lowered to
\texttt{arith.minf}/\texttt{arith.maxf}, which are not yet handled.

\textbf{Chained reductions.}
Expressions of the form \texttt{s = s + f(x) + g(y)} produce
intermediate SSA values between the load and the accumulating add,
breaking the load→op→store chain match in Phase~4.

\textbf{Non-constant loop bounds.}
When lower or upper bounds are not manifest constants, they are
reported as \texttt{?} in the output.  No analysis depends on them;
this is purely a reporting limitation.

\textbf{Confidence score heuristics.}
The confidence percentage displayed in the HTML report is computed in
\texttt{report.py} from the phase at which a verdict was reached and
the memory access pattern.  It is not derived from a formal
probability model and should be interpreted as a qualitative guide
only.

% =============================================================
\section{Future Work}
\label{sec:future}

\textbf{In-place update recognition.}
A use-def walk from the store operand back to the load can confirm that
the only value stored is a function of the loaded value at the same
address.  This would eliminate the false negative for \texttt{a(i) =
a(i) * k}.

\textbf{Multi-dimensional subscript analysis.}
Extending Phase~3 to analyze each subscript position independently
would allow outer-loop parallelism in perfectly nested loops
(\texttt{c(i,j) = ...}) to be confirmed when the outer subscript is
IV-derived regardless of the inner.

\textbf{Min/Max reduction support.}
Adding \texttt{arith.minf}, \texttt{arith.maxf}, \texttt{arith.mini},
and \texttt{arith.maxi} to the Phase~4 matcher and emitting the
corresponding \texttt{REDUCTION(MIN:)} / \texttt{REDUCTION(MAX:)}
clauses is straightforward.

\textbf{Interprocedural side-effect summaries.}
A module-level pre-pass could build a conservative side-effect summary
for each Fortran subroutine, tagging its dummy arguments as read-only,
write-only, or read-write.  Phase~2 could then consult this summary
when it encounters a call op, enabling parallelism recovery in loops
containing pure or elemental function calls.

\textbf{Polyhedral extension.}
Integrating the ISL polyhedral library~\cite{verdoolaege2010isl} via
MLIR's Presburger arithmetic support would allow affine dependence
tests on loops with non-unit step and multi-dimensional index
expressions, replacing the restricted IV/IV$\pm$k check in Phase~3.

\textbf{Source-level OpenMP rewriting.}
A companion transformation pass could automatically insert
\texttt{!\$OMP PARALLEL DO} directives into the original Fortran source
at the positions identified by FPA, with a dry-run mode that previews
changes before applying them.

\textbf{Integration with CI pipelines.}
Packaging FPA as a GitHub Action or pre-commit hook would allow HPC
teams to enforce that new Fortran commits do not introduce inadvertent
loop-carried dependencies.

% =============================================================
\section{Conclusion}
\label{sec:conclusion}

We have presented \textsc{FlangParallelAnalyzer}, a static analysis
tool that classifies Fortran \texttt{DO} loops for OpenMP
parallelization by operating on the FIR dialect of the LLVM Flang
compiler.  The five-phase analysis pipeline — structural metadata
collection, memory access classification, induction-variable index
pattern matching, scalar reduction detection, and conservative safety
fallback — is implemented as a standard MLIR pass and achieves 100\%
correct verdict classification on a 35-case test suite with zero false
positives.

The tool demonstrates that a relatively compact implementation
($\approx$\,700~LOC of C++ analysis code) built on mature compiler
infrastructure can automate a task that currently requires significant
manual effort from HPC developers.  By reporting the full
analysis trace alongside each verdict and exposing it through an
interactive HTML interface, FPA bridges the gap between opaque
compiler auto-parallelization and error-prone manual annotation.

FPA is available as open-source software and is ready for deployment
in any environment with the LLVM~18 toolchain.  The source code,
test suite, and build instructions are publicly available at
\url{https://github.com/AryanGupta21/FlangParallelAnalyzer}.

% =============================================================
\section*{Acknowledgment}

The authors thank the LLVM and Flang open-source communities for
their well-documented compiler infrastructure, which made this work
feasible within an academic project timeline.

% =============================================================
\begin{thebibliography}{00}

\bibitem{metcalf2011modern}
M.~Metcalf, J.~Reid, and M.~Cohen,
\textit{Modern Fortran Explained: Incorporating Fortran 2018},
Oxford University Press, 2018.

\bibitem{openmp51}
OpenMP Architecture Review Board,
``OpenMP Application Programming Interface, Version 5.1,''
\url{https://www.openmp.org/specifications/}, Nov. 2020.

\bibitem{allen1987automatic}
R.~Allen and K.~Kennedy,
``Automatic translation of {FORTRAN} programs to vector form,''
\textit{ACM Trans. Program. Lang. Syst.}, vol.~9, no.~4,
pp.~491--542, Oct. 1987.

\bibitem{wolfe1996high}
M.~Wolfe,
\textit{High Performance Compilers for Parallel Computing},
Addison-Wesley, 1996.

\bibitem{banerjee1988dependence}
U.~Banerjee,
\textit{Dependence Analysis for Supercomputing},
Kluwer Academic Publishers, 1988.

\bibitem{pugh1992omega}
W.~Pugh,
``The Omega test: a fast and practical integer programming algorithm for
dependence analysis,''
\textit{Commun. ACM}, vol.~35, no.~8, pp.~102--114, Aug. 1992.

\bibitem{blume1994polaris}
W.~Blume et~al.,
``Effective automatic parallelization with Polaris,''
\textit{Int. J. Parallel Program.}, vol.~23, no.~4, pp.~339--379, 1995.

\bibitem{hall1996maximizing}
M.~W.~Hall et~al.,
``Maximizing multiprocessor performance with the {SUIF} compiler,''
\textit{Computer}, vol.~29, no.~12, pp.~84--89, Dec. 1996.

\bibitem{quinlan2000rose}
D.~Quinlan,
``{ROSE}: Compiler support for object-oriented frameworks,''
\textit{Parallel Process. Lett.}, vol.~10, nos.~2--3, pp.~215--226, 2000.

\bibitem{llvm}
C.~Lattner and V.~Adve,
``{LLVM}: A compilation framework for lifelong program analysis and
transformation,''
in \textit{Proc. Int. Symp. Code Generation and Optimization (CGO)},
pp.~75--86, Mar. 2004.

\bibitem{lattner2021mlir}
C.~Lattner et~al.,
``{MLIR}: Scaling compiler infrastructure for domain specific computation,''
in \textit{Proc. IEEE/ACM Int. Symp. Code Generation and Optimization (CGO)},
pp.~2--14, 2021.

\bibitem{flang}
The LLVM Project,
``Flang: The {LLVM} {F}ortran compiler,''
\url{https://flang.llvm.org/}, 2024.

\bibitem{bondhugula2008practical}
U.~Bondhugula, A.~Hartono, J.~Ramanujam, and P.~Sadayappan,
``A practical automatic polyhedral parallelizer and locality optimizer,''
in \textit{Proc. ACM SIGPLAN Conf. Programming Language Design and
Implementation (PLDI)}, pp.~101--113, 2008.

\bibitem{parallelware}
Appentra Solutions,
``Parallelware Analyzer,''
\url{https://www.appentra.com/products/parallelware-analyzer/}, 2023.

\bibitem{verdoolaege2010isl}
S.~Verdoolaege,
``{isl}: An integer set library for the polyhedral model,''
in \textit{Proc. Int. Conf. Mathematical Software (ICMS)},
vol.~6327, pp.~299--302, 2010.

\end{thebibliography}

\end{document}
